\section{Xấp xỉ hàm cục bộ}
\label{sec:ch07-xap-xi-ham-cuc-bo}

\index{xấp xỉ hàm cục bộ}%
\index{hàm attribution}%
Ba mục vừa rồi cho thấy một bộ kỹ thuật đi qua cả ba nhánh, nhưng quan sát đó
chưa cho ai tính được gì. Khung hình thức duy nhất mà bài báo viện đến chỉ phủ
một nhánh: nó hợp nhất các phương pháp attribution đặc trưng dưới tên
\tn{xấp xỉ hàm cục bộ}{local function approximation}~\autocite{p18unified}.

Khung này đặt bài toán như sau. Quanh đầu vào $x$ có một lân cận $\mathcal{N}_x$
mà ta lấy mẫu nhiễu $\xi$ từ đó; có một lớp mô hình dễ đọc $\mathcal{G}$, thường
là các hàm tuyến tính; và có một hàm mất mát $\ell$ đo mức lệch giữa mô hình gốc
với mô hình dễ đọc trên các mẫu ấy. Lời giải thích là nghiệm của

\begin{equation}
\label{eq:ch07-lfa}
g^{*} = \operatorname*{arg\,min}_{g \in \mathcal{G}}\;
        \mathbb{E}_{\xi \sim \mathcal{N}_x}\,\ell(f, g, x, \xi).
\end{equation}

Bài báo gốc viết lân cận này là $\mathcal{Z}$; sách đổi sang $\mathcal{N}_x$ vì
chương~\ref{ch:lime} đã gắn $\mathcal{Z}$ với tập mẫu của LIME, và
phụ lục~\ref{app:ky-hieu} ghi lại phép đối chiếu đó. Chữ $g$ ở đây là mô hình
dễ đọc, và hệ số của nó chính là các điểm số mà hàm attribution $g$ của
định nghĩa~\ref{def:ch07-ham-attribution} gán, nên hai cách dùng cùng chỉ một
đối tượng.

Sức nặng của phương trình~\ref{eq:ch07-lfa} nằm ở số phương pháp mà nó chứa
được. Tám phương pháp attribution đặc trưng quen thuộc đều là trường hợp
riêng của nó, gồm Occlusion, KernelSHAP, Vanilla Gradients, Gradient$\mathbin{\times}$Input,
Integrated Gradients, SmoothGrad, LIME và C-LIME. Chúng dùng chung một lớp mô
hình tuyến tính $\mathcal{G}$ và chỉ khác nhau ở hai chỗ: lấy mẫu từ lân cận
nào, và đo sai lệch bằng hàm mất mát nào~\autocite{p18unified}. Một mẫu $\xi$
tác động lên $x$ theo một trong hai cách, nhân từng thành phần để che bớt đầu
vào hoặc cộng vào để làm nhiễu nó. Occlusion lấy
mẫu bằng vector one-hot ngẫu nhiên và đo bằng sai số bình phương; KernelSHAP
cũng đo bằng sai số bình phương nhưng lấy mẫu theo \tn{hàm nhân Shapley}{Shapley kernel};
LIME lấy mẫu theo một \tn{hàm nhân}{kernel} mũ; còn nhóm gradient thì đo bằng
sai lệch giữa gradient của $g$ và gradient của $f$ tại các mẫu, chứ không phải
sai số bình phương giữa hai đầu ra.

Điều kiện làm cho một hàm mất mát là hợp lệ cũng được nêu tường minh: kỳ vọng
của $\ell$ bằng không khi và chỉ khi $f$ và $g$ trùng nhau trên mọi mẫu rút từ
$\mathcal{N}_x$~\autocite{p18unified}. Nói cách khác, hàm mất mát phải phạt đúng
sự khác biệt giữa hai mô hình trong lân cận đã chọn, chứ không phạt một
\tn{đại lượng thay thế}{proxy} nào khác.

Như vậy khung xấp xỉ hàm cục bộ và khung độ đo của
chương~\ref{ch:attribution-tu-nguyen-ly-dau} là hai phép hợp nhất khác nhau cho
cùng một nhánh. Khung độ đo hỏi ta lấy trung bình đầu ra của mô hình theo trọng
số nào; khung xấp xỉ hàm cục bộ hỏi ta lấy mẫu ở đâu và đo sai lệch bằng gì. Cả
hai cùng đưa mọi phương pháp về một ký pháp, và cả hai cùng để nguyên lựa chọn
nằm ở giữa. Một phương pháp không sai vì chọn lân cận này thay vì lân cận kia;
nó chỉ đang trả lời một câu hỏi phản thực khác, và không khung nào trong hai
khung nói câu hỏi nào hợp với miền ứng dụng.

Với hai nhánh còn lại, bài báo chỉ nêu một giả thuyết, và giả thuyết ấy không
đối xứng. Ở nhánh đặc trưng, mô hình dễ đọc $g$ nhận một vector đặc trưng làm
đầu vào và xấp xỉ \emph{dự đoán} $f(x)$ của mô hình; đó là điều
phương trình~\ref{eq:ch07-lfa} viết. Các tác giả phỏng đoán rằng
\tn{attribution dữ liệu}{data attribution} xấp xỉ \emph{trọng số} của mô hình
trên không gian các điểm huấn luyện, còn
\tn{attribution thành phần}{component attribution} xấp xỉ dự đoán của mô hình
trên không gian các thành phần~\autocite{p18unified}. Sách diễn giải phỏng đoán ấy như sau: ở
nhánh dữ liệu, $g$ sẽ nhận một tập con của tập huấn luyện làm đầu vào và xấp
xỉ trọng số mà mô hình học được từ tập con ấy, tức là xấp xỉ kết quả của quá
trình huấn luyện chứ không phải một dự đoán; còn ở nhánh thành phần, $g$ sẽ
nhận một tập con các thành phần làm đầu vào và xấp xỉ dự đoán của mô hình khi
chỉ các thành phần ấy còn hoạt động. Các tác giả viết điều này ở dạng phỏng
đoán và không chứng minh, nên
\tn{hàm attribution}{attribution function} của
định nghĩa~\ref{def:ch07-ham-attribution} vẫn là mẫu số chung duy nhất mà cả
ba nhánh chắc chắn chia nhau.
